% ============================================================
%  Abstract
% ============================================================
\section*{Abstract}
\addcontentsline{toc}{section}{Abstract}

The International Maritime Organization's Carbon Intensity Indicator (\CII{})
framework, in force from 2023, assigns annual environmental ratings to vessels
based on the ratio of CO$_2$ emissions to a transport work proxy.  For
Roll-on/Roll-off (\RoRo{}) vessels, the denominator uses Gross Tonnage (\GT{})
rather than cargo weight, making speed and sailing distance the sole operational
levers available to an operator.

This thesis develops a Mixed-Integer Linear Programming (\MILP{}) formulation
for joint speed optimisation under \CII{} constraints, calibrated to the Baltic
Sea \RoRo{} service operated by SCA Logistics (M/V~Obbola, GT~19{,}918).
The model is validated against EU MRV operational data (153 voyage legs,
December 2021 -- August 2022), reproducing the audited attained \CII{} to within
1.2\,\%.  Six scenarios are analysed spanning current regulation (2026 Rating~C),
a stringent 2030 target (Rating~B), combined EU Emissions Trading System (\ETS{})
carbon pricing, and schedule flexibility.

The principal findings are threefold.  First, under the 2026 regulatory
framework, the published vessel timetable --- not the \CII{} target --- is the
binding constraint: the cost-optimal schedule yields an attained \CII{} of
10.27~g/\GT{}$\cdot$nm (Rating~A) with 33.9\,\% headroom below the C~boundary.
Speed optimisation tools provide limited near-term benefit; the regulatory
instrument is partially blunt for fixed-schedule short-sea services.  Second,
a 2030 Rating~B target (9.93~g/\GT{}$\cdot$nm, assuming a 31\,\% reduction
factor) is \emph{structurally infeasible} under the current timetable: four
bottleneck arcs with tight schedule windows force minimum speeds that together
produce an irreducible attained \CII{} of 10.32~g/\GT{}$\cdot$nm.  Achieving
Rating~B by 2030 requires schedule redesign, alternative fuels, or energy
efficiency investments --- speed management alone is insufficient.  Third,
replacing per-arc time budgets with a global cycle-time cap (scenario S5)
reduces per-loop fuel consumption by 15.2\,\% (134.7~t) and total voyage cost
by EUR~124{,}569 per loop (EUR~2.12~million per year), while further improving
the \CII{} rating to 8.71~g/\GT{}$\cdot$nm (Rating~A, 43.9\,\% slack).  The
monetised value of schedule flexibility is thus substantial and easily
justifiable against the cost of renegotiating port windows.

The thesis contributes a correctly specified \RoRo{} \CII{} formulation (using
\GT{}-based capacity and RoRo-specific rating multipliers from MEPC.339(76)),
a systematic infeasibility proof for the 2030 Rating~B target, and
scenario-based quantification of the economic and environmental trade-offs
facing Baltic Sea short-sea operators under evolving carbon regulation.

\medskip
\noindent\textbf{Keywords:} Carbon Intensity Indicator, maritime operations,
speed optimisation, short-sea shipping, mixed-integer linear programming,
Roll-on/Roll-off, Baltic Sea, EU ETS.
